% ============================================================
%  Colorimetric Analysis of Chemical Test Strips
%  Mathematical Methods for an Arsenic Colorimeter
%
%  Reconstructed from design conversation.
%  Compile with pdflatex (twice for cross-references).
%  Requires: amsmath, amssymb, booktabs, geometry, hyperref,
%            algorithm2e (optional), microtype.
% ============================================================

\documentclass[11pt, a4paper]{article}

\usepackage{amsmath, amssymb, amsthm}
\usepackage{booktabs}
\usepackage{geometry}
\geometry{margin=2.5cm}
\usepackage{hyperref}
\hypersetup{colorlinks=true, linkcolor=blue, citecolor=blue, urlcolor=blue}
\usepackage{microtype}
\usepackage{parskip}
\usepackage{enumitem}

% Convenience macros
\newcommand{\R}{\mathbb{R}}
\newcommand{\E}{\mathbb{E}}
\newcommand{\bx}{\mathbf{x}}
\newcommand{\bS}{\mathbf{S}}
\newcommand{\bI}{\mathbf{I}}
\newcommand{\bV}{\mathbf{V}}
\newcommand{\bC}{\mathbf{C}}
\newcommand{\bA}{\mathbf{A}}
\newcommand{\bPsi}{\boldsymbol{\Psi}}
\newcommand{\tr}{\operatorname{tr}}
\newcommand{\diag}{\operatorname{diag}}

\title{%
  Colorimetric Analysis of Chemical Test Strips:\\
  Mathematical Methods for a Bayer-Camera Colorimeter}
\author{Design Notes}
\date{}

\begin{document}
\maketitle
\tableofcontents
\newpage

% ============================================================
\section{Introduction and Instrument Overview}
% ============================================================

This document describes the mathematical methods developed for a
closed-box colorimeter designed to measure arsenic concentration in
water using chemical test strips.  The strips discolour in the range
white\,--\,yellow\,--\,brown, and the highest sensitivity is required
in the pale-yellow region where human colour discrimination is poor.

\subsection{Hardware}

\begin{itemize}[nosep]
  \item \textbf{Camera:} Sony IMX219, 8\,MP, controlled by a Raspberry Pi Zero.
    Raw Bayer images are captured; the four channels are
    $B$, $G_b$, $G_r$, $R$.
  \item \textbf{Illumination:} Narrow-band LEDs (blue, green, red, and
    optionally more), each used alone.  LED spectral output is
    temperature-dependent.
  \item \textbf{Stage:} A reference colour card with invariant patches
    (white, black, and several shades of yellow to brown) surrounds the
    test strip at the centre of the field.
\end{itemize}

\subsection{Key Design Principles}

\begin{enumerate}[nosep]
  \item Every image contains both the unknown test strip \emph{and} all
    invariant reference patches.  Color distances are computed
    \emph{within} a single image, making them independent of LED
    drive current and exposure time.
  \item LED temperature is retained as a continuous covariate because
    it shifts the LED spectrum in a way that is not cancelled by
    within-image normalisation.
  \item Calibration uses many test strips of known concentration spanning
    the measurement range.  Each strip produces a vector of colour
    distances; a Gaussian Process regression (GP) maps these to concentration.
\end{enumerate}

% ============================================================
\section{Image Channels and Normalisation}
% ============================================================

\subsection{Multi-spectral Channel Vector}

Each Bayer image under one LED colour yields four channels.
Only channels with sufficient signal-to-noise ratio (SNR) are used.
The SNR criterion for channel $k$ under LED $i$ is

\begin{equation}
  \mathrm{SNR}_{ik}
    = \frac{\bar{x}_{ik}^{\mathrm{white}}}{\sigma_{ik}^{\mathrm{dark}}}
    \geq \tau,
  \label{eq:snr}
\end{equation}

where $\bar{x}_{ik}^{\mathrm{white}}$ is the mean pixel value of the
white reference patch and $\sigma_{ik}^{\mathrm{dark}}$ is the
standard deviation of the dark (black patch) signal.
The threshold $\tau$ (typically 10) is fixed for all LEDs.
For example, the $R$ channel is near zero under blue illumination and
is excluded.

The set of channels passing the SNR test for LED $i$ is denoted
$\mathcal{C}_i \subseteq \{B, G_b, G_r, R\}$.
The total number of included channels across all LEDs is
$N = \sum_i |\mathcal{C}_i|$, and $N > 4$ in general.

\subsection{Normalisation}

Raw pixel means are normalised to remove illumination intensity and
exposure-time variation.
%For channel (color layer) $k$, patch (reference color area) $j$, LED (illumination color) $i$:
For color channel $k$, reference color patch $j$, LED illumination color $i$:

\begin{equation}
  \tilde{x}_{ijk}
    = \frac{\bar{x}_{ijk}}{t_{\exp,i}\;\bar{x}_{ik}^{\mathrm{white}}},
  \label{eq:normalise}
\end{equation}

where $t_{\exp,i}$ is the exposure time for LED $i$ and
$\bar{x}_{ik}^{\mathrm{white}}$ is from the \emph{same image}.
Both strip and reference patches are divided by the same white mean,
so the ratio is dimensionless and illumination-intensity-independent.
LED temperature drift is \emph{not} removed by this normalisation and
is handled explicitly as a regression covariate.

The normalised colour vector for a patch $j$ under LED $i$ is

\begin{equation}
  \tilde{\bx}_{ij}
    = \bigl(\tilde{x}_{ijk}\bigr)_{k \in \mathcal{C}_i}
    \in \R^{|\mathcal{C}_i|}.
\end{equation}

% ============================================================
\section{Covariance Estimation}
% ============================================================

\subsection{Pooled Within-Patch Covariance}

The Mahalanobis distance requires an estimate of the within-patch
pixel noise covariance $\bS$.  This is estimated once during
calibration by pooling per-pixel observations from all invariant
reference patches (excluding SAMPLE and BLACK) across all calibration
images taken with the same LED.

Let $\mathbf{x}_{ij\ell}$ denote the normalised $|\mathcal{C}_i|$-vector
of pixel $\ell$ in reference patch $j$ of a calibration image under
LED $i$, and let $n_{ij}$ be the pixel count of that patch.
The pooled sample covariance is

\begin{equation}
  \bS_i
    = \frac{1}{\displaystyle\Bigl(\sum_{j} n_{ij}\Bigr) - J}
      \sum_{j}\sum_{\ell=1}^{n_{ij}}
        \bigl(\mathbf{x}_{ij\ell} - \bar{\bx}_{ij}\bigr)
        \bigl(\mathbf{x}_{ij\ell} - \bar{\bx}_{ij}\bigr)^{\!\top},
  \label{eq:pooled-cov}
\end{equation}

where $J$ is the number of reference patches used.
Ledoit--Wolf shrinkage is applied to stabilise the estimate:

\begin{equation}
  \hat{\bS}_i = (1-\lambda)\,\bS_i + \lambda\,\hat{\sigma}^2\bI,
  \qquad
  \hat{\sigma}^2 = \frac{1}{p}\tr(\bS_i),
  \label{eq:lw}
\end{equation}

where $\lambda \in (0,1)$ is chosen analytically to minimise the
expected squared Frobenius error.  The result $\hat{\bS}_i$ is
positive definite with controllable condition number.

\subsection{Bayesian Update at Measurement Time}

When a new image is captured, the live reference patches provide
fresh pixel noise observations.  The calibration estimate is used as
an inverse-Wishart prior; the live scatter matrix updates it.

\paragraph{Prior.}
\begin{equation}
  \bS \sim \mathcal{W}^{-1}(\bPsi_0,\,\nu_0),
  \qquad
  \bPsi_0 = \nu_0\,\hat{\bS}_i,
  \label{eq:iw-prior}
\end{equation}

so that $\E[\bS] = \hat{\bS}_i$.
The prior weight $\nu_0$ is set to the total number of strided
calibration pixels used to estimate $\hat{\bS}_i$, reflecting the
confidence accumulated during calibration.

\paragraph{Live scatter matrix.}
From the $n_{\mathrm{live}}$ stride-subsampled pixels of the live
reference patches:

\begin{equation}
  \bA_{\mathrm{live}}
    = \sum_{j \in \mathcal{R}}\sum_{\ell=1}^{n_j}
        \bigl(\mathbf{x}_{j\ell} - \bar{\bx}_j\bigr)
        \bigl(\mathbf{x}_{j\ell} - \bar{\bx}_j\bigr)^{\!\top},
  \label{eq:live-scatter}
\end{equation}

where $\mathcal{R}$ denotes the set of invariant reference patches
and $n_{\mathrm{live}} = \sum_{j \in \mathcal{R}} n_j$.

\paragraph{Posterior mean.}
\begin{equation}
  \bS_{\mathrm{updated}}
    = \frac{\bPsi_0 + \bA_{\mathrm{live}}}
           {\nu_0 + n_{\mathrm{live}} - p - 1},
  \label{eq:iw-posterior}
\end{equation}

where $p = |\mathcal{C}_i|$ is the number of channels.
When $n_{\mathrm{live}} \ll \nu_0$ the posterior is dominated by the
calibration prior; when $n_{\mathrm{live}} \gg \nu_0$ the live data
dominates.  The update is most significant when LED temperature has
drifted from calibration conditions.

% ============================================================
\section{Mahalanobis Distance}
% ============================================================

\subsection{Definition}

The Mahalanobis distance between the normalised colour vector of the
test strip $\tilde{\bx}_{\mathrm{strip}}$ and that of reference patch
$r$, both in the same image under LED $i$, is

\begin{equation}
  d_M\!\left(\tilde{\bx}_{\mathrm{strip}},\,\tilde{\bx}_r\right)
    = \sqrt{
        \bigl(\tilde{\bx}_{\mathrm{strip}} - \tilde{\bx}_r\bigr)^{\!\top}
        \bS_i^{-1}
        \bigl(\tilde{\bx}_{\mathrm{strip}} - \tilde{\bx}_r\bigr)
      }.
  \label{eq:mahal}
\end{equation}

The squared distance $d_M^2$ is asymptotically $\chi^2_p$ distributed
under the hypothesis that the strip colour is drawn from the same
distribution as reference $r$.

\subsection{Why Within-Image Distances Are Stable}

Because both $\tilde{\bx}_{\mathrm{strip}}$ and $\tilde{\bx}_r$ are
normalised by the same white patch mean from the same image
(equation~\eqref{eq:normalise}), the distance is:
\begin{itemize}[nosep]
  \item independent of LED drive current and total flux;
  \item independent of camera exposure and gain;
  \item partially self-correcting for LED temperature drift (both
    patches shift together under a spectral change).
\end{itemize}
Residual temperature sensitivity arises from the interaction of the
LED spectral shift with the differential reflectance spectra of the
strip and the reference patches.  This is retained as a covariate
in the GP.

\subsection{Outlier Detection}

A distance is flagged as a statistical outlier if

\begin{equation}
  d_M^2 > \chi^2_{p,\,1-\alpha},
  \label{eq:chi2-flag}
\end{equation}

where $\chi^2_{p,\,1-\alpha}$ is the $(1-\alpha)$ quantile of the
chi-squared distribution with $p$ degrees of freedom.  At $\alpha = 0.01$
and $p = 4$ channels, the threshold is approximately $13.28$.
A flagged distance indicates a possible contaminated strip, wrong
reagent, or sensor artefact.

% ============================================================
\section{Feature Vector and Calibration Dataset}
% ============================================================

\subsection{Feature Vector}

For a test strip image set (one image per LED), the GP input feature
vector is assembled as follows.  Let $I$ denote the number of LEDs
and $\mathcal{R} = \{r_1, \ldots, r_J\}$ the set of reference patches.
The feature vector is

\begin{equation}
  \mathbf{f}
    = \Bigl(
        d_{11},\, d_{12},\, \ldots,\, d_{IJ},\;
        T_1,\, T_2,\, \ldots,\, T_I
      \Bigr)
    \in \R^{IJ + I},
  \label{eq:feature-vector}
\end{equation}

where $d_{ij} = d_M(\tilde{\bx}_{\mathrm{strip}},\, \tilde{\bx}_{r_j})$
is the Mahalanobis distance to reference $r_j$ under LED $i$, and
$T_i$ is the measured temperature of LED $i$ at capture time.

\subsection{Calibration Dataset}

Each calibration strip (known concentration $c$) yields one feature
vector $\mathbf{f}$ and one target $c$.  Repeating for $M$ strips:

\begin{equation}
  \mathcal{D}
    = \bigl\{(\mathbf{f}^{(m)},\, c^{(m)})\bigr\}_{m=1}^{M},
  \label{eq:cal-dataset}
\end{equation}

with $\mathbf{f}^{(m)} \in \R^{IJ+I}$ and $c^{(m)} \in \R$.
The calibration dataset spans the concentration range of interest,
with denser sampling in the pale-yellow region.

% ============================================================
\section{Gaussian Process Regression}
% ============================================================

\subsection{Model}

A Gaussian Process prior is placed over the function
$c = f(\mathbf{f})$:

\begin{equation}
  f \sim \mathcal{GP}\!\left(\mu(\mathbf{f}),\;
                              k(\mathbf{f}, \mathbf{f}')\right).
  \label{eq:gp-prior}
\end{equation}

The mean function $\mu$ is absorbed into a normalisation of the
target (zero-mean after standardisation).
The kernel is a product of two anisotropic squared-exponential
(RBF) kernels --- one for the distance features, one for the
temperature features --- plus a white noise term:

\begin{align}
  k(\mathbf{f}, \mathbf{f}')
    &= \sigma_f^2\;
       k_{\mathrm{dist}}(\mathbf{f}_d, \mathbf{f}_d')\;
       k_{\mathrm{temp}}(\mathbf{f}_T, \mathbf{f}_T')
       + \sigma_n^2\,\delta(\mathbf{f}, \mathbf{f}'),
  \label{eq:kernel} \\[4pt]
  k_{\mathrm{dist}}(\mathbf{f}_d, \mathbf{f}_d')
    &= \exp\!\left(
         -\frac{1}{2}
         \sum_{ij}
         \frac{(d_{ij} - d_{ij}')^2}{\ell_{ij}^2}
       \right),
  \label{eq:kernel-dist} \\[4pt]
  k_{\mathrm{temp}}(\mathbf{f}_T, \mathbf{f}_T')
    &= \exp\!\left(
         -\frac{1}{2}
         \sum_{i}
         \frac{(T_i - T_i')^2}{\ell_{T,i}^2}
       \right).
  \label{eq:kernel-temp}
\end{align}

Here $\mathbf{f}_d$ and $\mathbf{f}_T$ are the distance and temperature
sub-vectors of $\mathbf{f}$; $\sigma_f^2$ is the signal variance;
$\sigma_n^2$ is the observation noise variance; and $\ell_{ij}$,
$\ell_{T,i}$ are per-feature length scales.
All hyperparameters are optimised by maximising the log marginal
likelihood:

\begin{equation}
  \log p(\mathbf{c} \mid \mathbf{F})
    = -\frac{1}{2}\mathbf{c}^{\top}
        \bigl(\mathbf{K} + \sigma_n^2\bI\bigr)^{-1}\mathbf{c}
      -\frac{1}{2}\log\det\!\bigl(\mathbf{K} + \sigma_n^2\bI\bigr)
      -\frac{M}{2}\log 2\pi,
  \label{eq:log-ml}
\end{equation}

where $\mathbf{K}_{mm'} = k(\mathbf{f}^{(m)}, \mathbf{f}^{(m')})$
is the $M \times M$ kernel matrix and
$\mathbf{c} = (c^{(1)}, \ldots, c^{(M)})^{\top}$.

\subsection{Prediction}

Given a new feature vector $\mathbf{f}^*$, the GP posterior
predictive distribution is Gaussian with mean and variance

\begin{align}
  \hat{c}
    &= \mathbf{k}_*^{\top}
       \bigl(\mathbf{K} + \sigma_n^2\bI\bigr)^{-1}\mathbf{c},
  \label{eq:gp-mean} \\[4pt]
  \hat{\sigma}^2
    &= k(\mathbf{f}^*, \mathbf{f}^*)
       - \mathbf{k}_*^{\top}
         \bigl(\mathbf{K} + \sigma_n^2\bI\bigr)^{-1}
         \mathbf{k}_*,
  \label{eq:gp-var}
\end{align}

where $(\mathbf{k}_*)_m = k(\mathbf{f}^*, \mathbf{f}^{(m)})$.
The estimate $\hat{c}$ is the predicted concentration and
$\hat{\sigma}$ is the corresponding $1\sigma$ uncertainty.

A measurement is flagged as potentially extrapolated if
$\mathbf{f}^*$ lies outside the per-feature range of the calibration
set $\mathcal{D}$.

% ============================================================
\section{PCA Projection (Supplementary)}
% ============================================================

When the colour gamut varies along essentially one physical axis
(white\,--\,yellow\,--\,brown driven by a single chemical), the
between-reference covariance has one dominant eigenvalue and a 1-D
projection suffices for a scalar calibration curve.

Let the $N_c$ normalised reference-patch colour vectors
$\{\bar{\bx}_r\}_{r \in \mathcal{R}}$ have centroid
$\bar{\bx} = \frac{1}{|\mathcal{R}|}\sum_r \bar{\bx}_r$.
The between-reference covariance is

\begin{equation}
  \bC = \frac{1}{|\mathcal{R}|-1}
        \sum_{r \in \mathcal{R}}
        (\bar{\bx}_r - \bar{\bx})(\bar{\bx}_r - \bar{\bx})^{\top}.
  \label{eq:between-cov}
\end{equation}

Its eigendecomposition $\bC\mathbf{v}_k = \lambda_k \mathbf{v}_k$
yields eigenvectors sorted in descending order.
The scalar projection of a colour vector onto the principal axis is

\begin{equation}
  s = \mathbf{v}_1^{\top}(\bx - \bar{\bx}).
  \label{eq:pca-project}
\end{equation}

The explained-variance ratio $\lambda_1 / \sum_k \lambda_k$ quantifies
how well the 1-D approximation holds.  If this ratio is close to 1,
the GP distance-based approach and the PCA scalar approach will give
nearly identical results; the GP approach is more general and is
preferred.

% ============================================================
\section{Summary of Measurement Procedure}
% ============================================================

\begin{enumerate}
  \item For each available LED $i$, capture a Bayer image of the
    test strip on the stage.

  \item For each responsive channel $k \in \mathcal{C}_i$, extract
    the patch means for SAMPLE and all reference patches.

  \item Optionally update the calibration covariance using
    equation~\eqref{eq:iw-posterior} with the live reference pixels.

  \item Compute the normalised colour vectors
    (equation~\eqref{eq:normalise}) for the SAMPLE and each
    reference patch.

  \item Compute the Mahalanobis distances
    $d_{ij}$ (equation~\eqref{eq:mahal}) and check for outliers
    (equation~\eqref{eq:chi2-flag}).

  \item Read the LED temperature $T_i$ for each LED.

  \item Assemble the feature vector $\mathbf{f}$
    (equation~\eqref{eq:feature-vector}).

  \item Query the fitted GP to obtain
    $(\hat{c}, \hat{\sigma})$ (equations \eqref{eq:gp-mean}
    and~\eqref{eq:gp-var}).

  \item Report $\hat{c}$ as the concentration estimate with
    $\hat{\sigma}$ as the $1\sigma$ uncertainty.
    Raise warnings for chi-squared outliers or extrapolation.
\end{enumerate}

% ============================================================
\section*{References}
% ============================================================

\begin{enumerate}
  \item Mahalanobis, P.\,C.\ (1936).
    On the generalised distance in statistics.
    \textit{Proceedings of the National Institute of Sciences of India},
    \textbf{2}(1), 49--55.

  \item Duda, R.\,O., Hart, P.\,E., \& Stork, D.\,G.\ (2001).
    \textit{Pattern Classification} (2nd ed.).
    Wiley-Interscience.
    Chapter~2: Bayesian decision theory and Mahalanobis distance.

  \item McLachlan, G.\,J.\ (1999).
    Mahalanobis distance.
    \textit{Resonance}, \textbf{4}(6), 20--26.

  \item Rasmussen, C.\,E.\ \& Williams, C.\,K.\,I.\ (2006).
    \textit{Gaussian Processes for Machine Learning}.
    MIT Press.
    \url{http://www.gaussianprocess.org/gpml/}

  \item Geladi, P., Grahn, H., \& Burger, J.\ (2007).
    Multivariate images, hyperspectral imaging.
    In \textit{Techniques and Applications of Hyperspectral Image
    Analysis}. Wiley.
    Covers multi-illumination image stacks as joint multivariate
    observations.

  \item Ledoit, O.\ \& Wolf, M.\ (2004).
    A well-conditioned estimator for large-dimensional covariance matrices.
    \textit{Journal of Multivariate Analysis}, \textbf{88}(2), 365--411.
\end{enumerate}

\end{document}